\subsection{CENTRAL FILTRATIONS ON A GROUP}

DEFINITION 2. Let G be a group. A real filtration \((\mathrm{G}_{\alpha})\) on G is called central if \(\mathrm{G} = \bigcup_{\alpha >0}\mathrm{G}_{\alpha}\) and the commutator \((x,y) = x^{-1}y^{-1}xy\) of an element \(x\) of \(\mathrm{G}_{\alpha}\) and an element \(y\) of \(\mathrm{G}_{\beta}\) belongs to \(\mathrm{G}_{\alpha +\beta}\) .

In terms of the order function v, the above definition translates into the relations

\[
v (x) > 0, \quad v ((x, y)) \geqslant v (x) + v (y) \quad \text { for   all } x, y \text { in } G.\tag{10}
\]

We deduce that \(v((x,y)) > v(x)\) if \(v(x) \neq +\infty\) ; if we write \(x^y = y^{-1}xy\) (cf.~Algebra, Chapter I, § 6, no. 2), then \(x^y = x.(x,y)\) , whence

\[
v (x ^ {y}) = v (x).\tag{11}
\]

This relation expresses the fact that each of the subgroups \(G_{\alpha}\) of G is normal. The \(G_{\alpha}\) form a fundamental system of neighbourhoods of e for a topology compatible with the group structure on G (General Topology, Chapter III, § 1, no. 2, Example) said to be defined by the filtration ( \(G_{\alpha}\) ).

In the rest of this no., G denotes a group with a central filtration \((G_{\alpha})\) . For all \(\alpha \in R\) , we define the subgroup \(G_{\alpha}^{+}\) of G by

\[
\mathrm{G} _ {\alpha} ^ {+} = \bigcup_ {\beta > \alpha} \mathrm{G} _ {\beta}.\tag{12}
\]

In particular \(G_{\alpha}^{+} = G_{\alpha} = G\) for \(\alpha \leqslant 0\) . Recall that if A and B are two subgroups of G, (A, B) denotes the subgroup of G generated by the commutators \((a, b)\) with \(a \in A\) and \(b \in B\) . With this notation we have the formulae

\begin{enumerate}
\def\labelenumi{(\arabic{enumi})}
\setcounter{enumi}{12}
\tightlist
\item
\end{enumerate}

\[
\left(\mathrm{G} _ {\alpha}, \mathrm{G} _ {\beta}\right) \subset \mathrm{G} _ {\alpha + \beta}\tag{13}
\]

\[
\left(\mathrm{G} _ {\alpha} ^ {+}, \mathrm{G} _ {\beta}\right) \subset \mathrm{G} _ {\alpha + \beta} ^ {+}\tag{13'}
\]

\[
(G, G _ {\alpha}) \subset G _ {\alpha} ^ {+}.\tag{14}
\]

By (14), \(\mathrm{G}_{\alpha}^{+}\) is a normal subgroup of \(\mathrm{G}_{\alpha}\) for all \(\alpha \in \mathbf{R}\) and the quotient group \(\mathrm{gr}_{\alpha}(\mathrm{G}) = \mathrm{G}_{\alpha} / \mathrm{G}_{\alpha}^{+}\) is commutative. We write \(\mathrm{gr}(\mathrm{G}) = \bigoplus_{\alpha \in \mathbb{R}}\mathrm{gr}_{\alpha}(\mathrm{G})\) and give this group the graduation of type \(\mathbf{R}\) in which \(\mathrm{gr}_{\alpha}(\mathrm{G})\) consists of elements of degree \(\alpha\) . Then \(\mathrm{gr}_{\alpha}(\mathrm{G}) = \{0\}\) for \(\alpha \leqslant 0\) .

PROPOSITION 1. (i) Let \(\alpha, \beta\) be in \(\mathbf{R}\) . There exists a biadditive mapping

\[
\phi_ {\alpha \beta}: \operatorname{gr} _ {\alpha} (\mathrm{G}) \times \operatorname{gr} _ {\beta} (\mathrm{G}) \rightarrow \operatorname{gr} _ {\alpha + \beta} (\mathrm{G})
\]

which maps \((x\mathrm{G}_{\alpha}^{+},y\mathrm{G}_{\beta}^{+})\) onto \((x,y)\mathrm{G}_{\alpha +\beta}^{+}\)

\begin{enumerate}
\def\labelenumi{(\roman{enumi})}
\setcounter{enumi}{1}
\item
  Let \(\phi\) be the biadditive mapping of \(\operatorname{gr}(\mathbf{G}) \times \operatorname{gr}(\mathbf{G})\) into \(\operatorname{gr}(\mathbf{G})\) whose restriction to \(\operatorname{gr}_{\alpha}(\mathbf{G}) \times \operatorname{gr}_{\beta}(\mathbf{G})\) is \(\phi_{\alpha \beta}\) for every ordered pair \((\alpha, \beta)\) . The mapping \(\phi\) gives \(\operatorname{gr}(\mathbf{G})\) a Lie Z-algebra structure.
\item
  Recall the identity
\end{enumerate}

\[
(x x ^ {\prime}, y) = (x, y) ^ {x ^ {\prime}} \left(x ^ {\prime}, y\right)\tag{15}
\]

for \(x, x', y\) in \(\mathbf{G}\) (Algebra, Chapter I, § 6, no. 2, formula (4 bis)).

For \(x \in G_{\alpha}\) and \(y \in G_{\beta}\) , the class modulo \(G_{\alpha + \beta}^{+}\) of the element \((x, y)\) of \(G_{\alpha + \beta}\) will be denoted by \(f(x, y)\) . For \(a\) in \(G_{\alpha + \beta}\) and \(x'\) in \(G, a^{-1}.a^{x'} = (a, x') \in G_{\alpha + \beta}^{+}\) ; in particular \(f(x, y)\) is equal to the class modulo \(G_{\alpha + \beta}^{+}\) of \((x, y)^{x'}\) . Formula (15) therefore implies

\[
f (x x ^ {\prime}, y) = f (x, y) f (x ^ {\prime}, y).\tag{16}
\]

Now \((y, x) = (x, y)^{-1}\) , whence

\[
f (y, x) = f (x, y) ^ {- 1}.\tag{17}
\]

From (16) and (17) we deduce

\[
f (x, y y ^ {\prime}) = f (x, y) f (x, y ^ {\prime}).\tag{18}
\]

We have to prove that the mapping \(f\colon \mathrm{G}_{\alpha}\times \mathrm{G}_{\beta}\to \mathrm{gr}_{\alpha +\beta}(\mathrm{G})\) defines on taking quotients a mapping \(\phi_{\alpha \beta}\colon \mathrm{gr}_{\alpha}(\mathrm{G})\times \mathrm{gr}_{\beta}(\mathrm{G})\to \mathrm{gr}_{\alpha +\beta}(\mathrm{G})\) . By (16) and (18) it suffices to prove that \(f(x,y) = 0\) if \(x\in \mathrm{G}_{\alpha}^{+}\) or \(y\in \mathrm{G}_{\beta}^{+}\) , which follows from (13').

\begin{enumerate}
\def\labelenumi{(\roman{enumi})}
\setcounter{enumi}{1}
\tightlist
\item
  As \((x, x) = e\) , it follows from (17) that \(\phi\) is an alternating \(\mathbf{Z}\) -bilinear mapping. Hence it remains to prove that, for \(u \in \mathrm{gr}_{\alpha}(\mathrm{G})\) , \(v \in \mathrm{gr}_{\beta}(\mathrm{G})\) and \(w \in \mathrm{gr}_{\gamma}(\mathrm{G})\) ,
\end{enumerate}

\[
\phi (u, \phi (v, w)) + \phi (v, \phi (w, u)) + \phi (w, \phi (u, v)) = 0.\tag{19}
\]

Let \(x \in G_{\alpha}, y \in G_{\beta}\) and \(z \in G_{\gamma}\) be elements representing respectively \(u, v\) and \(w\) . We know that \(x^{y}\) and \(x\) are two elements of \(G_{\alpha}\) which are congruent modulo \(G_{\alpha}^{+}\) and hence \(x^{y}\) is a representative of \(u\) in \(G_{\alpha}\) ; as \((y, z)\) is a representative of \(\phi(v, w)\) in \(G_{\beta + \gamma}\) , we see that \((x^{y}, (y, z))\) is a representative of \(\phi(u, \phi(v, w))\) in \(G_{\alpha + \beta + \gamma}\) . By cyclic permutation, we see that \((y^{z}, (z, x))\) and \((z^{x}, (x, y))\) represent respectively \(\phi(v, \phi(w, u))\) and \(\phi(w, \phi(u, v))\) in \(G_{\alpha + \beta + \gamma}\) . Relation (19) is then a consequence of the following identity (Algebra, Chapter I, §6, no. 2, formula (15)):

\[
\left(x ^ {y}, (y, z)\right). \left(y ^ {z}, (z, x)\right). \left(z ^ {x}, (x, y)\right) = e.\tag{20}
\]

The Lie algebra \(\mathrm{gr}(\mathbf{G})\) over \(\mathbf{Z}\) defined in Proposition 1 is called the graded Lie algebra associated with the filtered group \(\mathbf{G}\) .

